\runninghead{Both Grow until the Harvest}
\properlane{patience-of-the-field}{Both Grow until the Harvest: Judgment Left to the Reapers, Forbearance Asked of the Servants}

The householder refuses his servants' offer to weed the field and gives his
reason in the parable itself: \latin{ne forte colligéntes zizánia,
eradicétis simul cum eis et tríticum}, ``lest perhaps gathering up the cockle,
you root up the wheat also together with it.'' The Fathers who comment on it
find two facts about the present age behind that reason: the wheat is in
danger from those who would weed, and the cockle may yet become wheat. St
Augustine gives the servants' own fallibility as the reason they were
refused, and tells the tares in a separate address that they may still
become wheat. St Jerome finds the danger in the plant, which in the blade
cannot be told from wheat, and reads room for repentance in \latin{ne forte}.
St John Chrysostom reads both facts in that clause, and takes the command as
a prohibition of killing. The separation of persons then belongs to the
harvest and to the angels whom the Lord names as its reapers. The servants
still have work before it, and this interpretation finds it in Paul's
requests to the Colossians: patience, bearing with one another, and forgiving
as the Lord forgave.

\subsection{The servants' question and the Apostle's command}
\label{sec:patience-question}

The servants see weeds among the wheat,
they know their master sowed good seed, and they offer to do what any farmer's
men would do. The New Testament also seems to be on their side. St Paul tells
the Corinthians to expel the man living in incest, and St Jerome quotes him
against the parable: \latin{Sinite utraque crescere} appears to contradict
\latin{Auferte malum de medio vestrum}, ``take away the evil from among
you'' (1~Cor 5:13), and if the Church is commanded both to remove the wicked
and to leave them, it needs to know which command applies when (\work{In
Matthaeum} II, on 13:36--43, PL~26 cols.~93--94). St Thomas Aquinas praises
the servants on two counts: they are moved by zeal to destroy evil, as the
Apostle commands, and they will not do it of their own motion but wait for
their lord's order, \latin{iussu domini} (\work{Super Mt.}\ c.~13). St John
Chrysostom says the same of their haste: it ``shows their anxiety for the
crop,'' and they ``trust not themselves with it, but await the Master's
decision, saying, `Wilt Thou?''' (\work{Hom.\ in Mt.}\ 46.1).

The question is therefore what the householder's refusal forbids, and on
what ground.

\subsection{St Jerome: the blade that cannot be told apart}
\label{sec:patience-jerome}

St Jerome answers from the plant. Between wheat and \latin{lolium}, darnel,
there is great likeness while they are still in the blade, so that the two
can hardly be distinguished, or not at all. The Lord's
command, he says, teaches \latin{ne ubi quid ambiguum est, cito sententiam
proferamus; sed Deo judicii terminum reservemus}: ``that where anything is
doubtful we should not quickly pass sentence, but reserve to God the end of
judgment,'' so that what God removes at the end is \latin{non suspicionem
criminis, sed manifestum reatum}, ``not a suspicion of crime, but manifest
guilt'' (PL~26 cols.~93--94).

His second point comes from \latin{ne forte \dots\ eradicetis}: \latin{datur
locus poenitentiae, et monemur ne cito amputemus fratrem: quia fieri potest,
ut ille qui hodie noxio depravatus est dogmate, cras resipiscat}, ``room is
given for repentance, and we are warned not to cut off a brother quickly,
because it can happen that one who today is corrupted by a harmful teaching
will tomorrow come to his senses.''

\subsection{St Augustine: ``ye are such as may be deceived''}
\label{sec:patience-augustine}

St Augustine preached on this Gospel the day after preaching on the sower,
and in his sermon he speaks to three groups in turn (\work{Sermo}~73). To the
zealous servants he speaks in the Lord's voice: ``Why are ye so hasty, He
says, ye servants full of zeal? Ye see tares among the wheat, ye see evil Christians among the good; and
ye wish to root up the evil ones; be quiet, it is not the time of harvest.
That time will come, may it only find you wheat!'' And: ``In the field they
may be with you, but they will not be so in the barn'' (73.1).

Then, in his own voice, ``to-day I am addressing the tares,'' whom he
identifies with the wayside, the stony ground and the thorns of the sower and
calls evil Christians: ``amend yourselves before the harvest come \dots\ it may so be, that they who to-day are tares, may to-morrow
be wheat'' (73.3).

``And so I will address the wheat also.'' To those whose lives are good and
who ``sigh and groan as being few among many'' he says, ``The angels will
come who can make the separation, and who cannot make mistakes.'' Then he
places himself and his hearers among the servants: ``We in this time present
are like those servants of whom it was said, `Wilt Thou that we go and gather
them up?' for we were wishing, if it might be so, that no evil ones should
remain among the good.'' They were refused, he says, ``For ye are such as may
be deceived,'' and ``We are but men, the reapers are the angels'' (73.4).

Augustine then asks his hearers whether they think the tares are ``all below,
and none above up here,'' and answers: ``even in these high seats there is
both wheat, and tares, and among the laity there is wheat, and tares.'' The
``high seats'' are the bench in the apse where the bishop sat with his
presbyters. His charge to both is one sentence: ``Let the good tolerate the
bad; let the bad change themselves, and imitate the good.''

In his expositions of the Psalms Augustine describes the state of the wheat
during the waiting. In this pilgrimage, he says, ``I am to live together with
the `tares' even to the time of the `harvest'''; the present is ``a certain
distinction without separation,'' the tares being still tares and the corn
still corn, and patience is needed ``to endure, until the harvest''
(\work{Enarr.\ in Ps.}\ 42.2--3). And living among the tares does the wheat no
harm: ``The wheat is not choked by the tares, nay, through endurance of the
tares the fruit of the wheat is increased'' (\work{Enarr.\ in Ps.}\ 64.15). In
his questions on the Gospels he assigns the separation to the good angels,
since \latin{officia vindictae possunt implere boni bono animo}, the good can
carry out the offices of punishment with a good mind (\work{Quaest.\ Ev.}\
I.10).

\subsection{St John Chrysostom: the prohibition of the sword}
\label{sec:patience-chrysostom}

St John Chrysostom takes the cockle to be heretics, and asks why the servants
are brought into the story at all: ``That He may pronounce it wrong to slay
them.'' The Lord, he says, forbids the uprooting ``to hinder wars from arising,
and blood and slaughter,'' and gives two reasons in the words \latin{ne forte
\dots\ eradicetis}: ``If ye are to take up arms, and to kill the heretics,
many of the saints also must needs be overthrown with them; or that of the
very tares it is likely that many may change and become wheat. If therefore ye
root them up beforehand, ye injure that which is to become wheat, slaying
some, in whom there is yet room for change and improvement'' (\work{Hom.\ in
Mt.}\ 46.1).

The prohibition has a limit: ``He doth not therefore forbid our checking
heretics, and stopping their mouths, and taking away their freedom of speech,
and breaking up their assemblies and confederacies, but our killing and
slaying them.'' And it has an end. Of the harvest Chrysostom says, ``So long as
they stand by the wheat, we must spare them, for it is possible for them even
to become wheat but when they have departed, having profited nothing, then of
necessity the inexorable punishment will overtake them.''

\subsection{How far correction may go}
\label{sec:patience-discipline}

Jerome, Augustine and Chrysostom do not agree on how much correction the
command allows now. Jerome reserves the doubtful case to God and does not
exclude putting away the manifest one. Chrysostom forbids killing and allows
restraint short of it.

St Augustine's own position is in his letter to Vincentius, a Donatist
bishop. ``Originally my opinion was,'' he writes, ``that no one should be
coerced into the unity of Christ, that we must act only by words, fight only
by arguments, and prevail by force of reason, lest we should have those whom
we knew as avowed heretics feigning themselves to be Catholics.'' He gave it
up, he says, not because of his opponents' arguments but because of examples,
his own town first, which had been brought into Catholic unity ``by fear of the
imperial edicts'' (\work{Ep.}~93.17). The penalties he defends are ``such as
exile and fines,'' which he calls ``moderate severity, or rather clemency,''
aimed at admonishing those who suffer them rather than punishing them; he
contrasts them with the death appointed by imperial law for pagan sacrifice
(93.10). In the same letter the parable forbids separation before the harvest.
The Church, sown by the Son of man, grows among the tares until the end, ``for
the field is the world, and the harvest is the end of time'' (93.31); she
looks for the visible separation ``only on the sea-shore, i.e.\ at the end of
the world,--meanwhile correcting as many as she can, and bearing with those
whom she cannot correct'' (93.34); and St Cyprian rebuked those who would
separate themselves as righteous, ``assuming to themselves that which the
Lord did not concede even to apostles,--namely, the gathering of the tares
before the harvest'' (93.36). Augustine argues for the penalties from their
effects, not from the parable, and his sermon on the parable names no
penalty.

St Thomas Aquinas, commenting on Matthew, gives four reasons the wicked are
not to be uprooted: the good are exercised by them; one who is wicked now may
become good, \latin{ut Paulus}, as Paul did; some seem wicked and are not, so
\latin{Nolite ante tempus iudicare}, ``judge not before the time'' (1~Cor
4:5); and a powerful man drags many down with him (\work{Super Mt.}\ c.~13).
He reports Chrysostom's reading and Augustine's change of mind, and holds the
prohibition where uprooting cannot be done without danger, \latin{si non
possit sine periculo}.

In the \work{Summa} Aquinas reports the same two positions: some have
understood the command to forbid ``not the excommunication but the slaying of
heretics, as appears from the words of Chrysostom,'' and Augustine says in his
letter to Vincentius that he changed his mind. He then takes the conditional
reading from Augustine: the meaning of ``Suffer both to grow until the
harvest'' ``must be gathered from those which precede, `lest perhaps gathering
up the cockle, you root the wheat also together with it,''' and he cites
Augustine's work against Parmenianus (III.2) for the rule that when this is
not to be feared, ``when a man's crime is so publicly known, and so hateful to
all, that he has no defenders, or none such as might cause a schism, the
severity of discipline should not slacken'' (II-II q.~10 a.~8 ad~1).

In the next question he concludes that heretics deserve not only
excommunication but death. On their side he gives the sin, ``for it is a much
graver matter to corrupt the faith which quickens the soul, than to forge
money, which supports temporal life.'' On the Church's side he gives mercy,
which condemns ``not at once, but `after the first and second admonition,'''
and then, if the heretic is still stubborn, ``no longer hoping for his
conversion, looks to the salvation of others'': she excommunicates him ``and
furthermore delivers him to the secular tribunal to be exterminated thereby
from the world by death,'' and Aquinas adds St Jerome on Gal 5:9, ``Cut off
the decayed flesh'' (q.~11 a.~3). He meets the parable in answering an
objection, with the reading he has already given: the Lord's command ``is to
be understood as referring to the case when the cockle cannot be plucked up
without plucking up the wheat'' (ad~3). The letter to Vincentius, which he
cites, accepts exile and fines for correction and does not ask for death.

\subsection{St Rabanus Maurus: the patience of one who is wronged}
\label{sec:patience-rabanus}

St Rabanus Maurus adds a moral exposition that turns from the Church's
judgment of others to each man's response to injury. The householder knew it
was an enemy's doing, and yet he did not blaze up in anger or seek to avenge
himself at once, but waited. \latin{Ita et nos cum
aliquam laesionem ab inimicis nostris acceperimus, non statim ad vindictam
prorumpere, sed magis ad virtutem patientiae recurrere debemus}: ``so we too,
when we have received some injury from our enemies, ought not to break out at
once into revenge, but rather to have recourse to the virtue of patience.''
\latin{Aequanimiter ergo debemus sufferre malos propter bonos}, ``we ought to
bear the wicked with an even mind for the sake of the good,'' lest, provoked
by them and repaying them, we give the good an occasion of scandal. The
command \latin{Sinite} builds longanimity, which must be kept not for a short
time but \latin{usque ad finem vitae}, to the end of life, because, as the
Lord says, he who perseveres to the end will be saved (Mt 24:13).

At the harvest Rabanus turns to what the householder does then:
\latin{Qui ante per longanimitatem patienter sustinuit, modo per justitiam
meritis condigna rependit}, ``he who before bore patiently through
long-suffering now repays by justice what the deserts merit.'' The same
householder who waited is the one who judges, and in that, Rabanus says, he
gives us an example, \latin{ut in patiendo justitiam meditemur, et in
judicando patientiam habere non negligamus}: ``that in suffering we may keep
justice in mind, and in judging not neglect to have patience'' (\work{In
Matthaeum} IV, on 13:24--30, PL~107 cols.~947--948).

\subsection{Bearing with one another}
\label{sec:patience-epistle}

Paul's list ends in \latin{patiéntiam}, and his next words say how patience
is to act:
\latin{supportántes ínvicem, et donántes vobismetípsis, si quis advérsus
áliquem habet querélam: sicut et Dóminus donávit vobis, ita et vos}. St John
Chrysostom explains \latin{supportantes} as ``passing things over,'' and the
pattern ``even as Christ forgave you'' as an obligation to forgive serious
wrongs and to love the offender, ``even if we have not been the first to
injure'' (\work{Hom.\ in Col.}\ 8). St Thomas Aquinas explains why bearing
with one another is needed at all: a man may keep to justice himself and still
find that \latin{aliorum mores sunt ei importabiles}, the ways of others are
unbearable to him, as Lot was tormented living among the Sodomites (2~Pet
2:8). One forgives, Aquinas says, when one holds no rancour and procures no
harm; \latin{sed quando necessitas puniendi est, tunc puniendus est}, ``but
when there is need to punish, then he is to be punished'' (\work{Super Col.}\
c.~3 lect.~3).

Rabanus, on the parable, asks us to bear injury without revenge, and Aquinas,
on the Epistle, asks a man to bear the ways of others that he finds
unbearable; neither mentions the other's text. Both describe the same act,
bearing with others instead of repaying or removing them, and it is the act
Paul names in \latin{supportántes ínvicem}, ``bearing with one another.''

The Introit comes from the letter in which Jeremias tells the exiles to seek
the peace of the city that holds them captive and to pray to the Lord for it
(Jer 29:7). St Jerome reads that verse as a command to the Church: we are to
pray for our enemies and seek their salvation as well as our own
(\work{In Hieremiam} V, on 29:7, PL~24 cols.~891--892). The antiphon sings
the promise the same letter makes to those who pray, \latin{invocábitis me,
et ego exáudiam vos}, ``you shall call upon me, and I will hear you.'' Since
the command and the promise stand in one letter, the promise of a hearing
covers the prayer for the city that the letter commands. Servants who may not
weed the field can still make that prayer for the cockle, asking for the
change Augustine's sermon asks of the tares.

The Communion sings the Lord's promise that prayer made in faith receives
what it asks, \latin{quidquid orántes pétitis, crédite quia accipiétis}. The
next verse in Mark, which the Mass does not sing, tells
those who pray to forgive first, and St Bede reads the two together: those
not yet perfect obtain what they ask ``provided that they first forgive
others who sin against them.'' The forgiveness Bede makes the condition of
being heard is the act Paul asks of the Colossians in \latin{donántes
vobismetípsis}.

\subsection{Who frees the field}
\label{sec:patience-who-frees}

The Gradual thanks God: \latin{Liberásti nos, Dómine, ex affligéntibus nos: et eos, qui
nos odérunt, confudísti}. St Thomas Aquinas, expounding the psalm, says the
saints are afflicted in two ways, in body and \latin{spiritualiter, quando
peccant illi qui habitant apud eos}, ``spiritually, when those who live among
them sin,'' \latin{et ab his afflictionibus liberantur a Deo}, ``and from
these afflictions they are freed by God'' (\work{Super Ps.}\ 43 n.~3).
His second kind of affliction describes the wheat's situation in the parable,
which lives beside the cockle. The Collect asks God to keep his family with
continual mercy and to defend it always by his protection, since it
\latin{in sola spe grátiæ cæléstis innítitur}, leans on the hope of heavenly
grace alone. The servants' offer to weed came, Chrysostom says, from anxiety
for the crop. Since they may not weed, that care is left to this petition,
made to God, who in Aquinas's sentence is the one who frees.

The Alleluia and the Offertory cry from the depths. ``Who is it then who
crieth?'' St Augustine asks of Psalm~129, and answers, ``A sinner''
(\work{Enarr.\ in Ps.}\ 129.2). If that is so, those who sing the cry and
would weed the field are themselves among those who need mercy, and in the
Secret that follows they ask God to absolve \latin{delícta nostra}, their own
sins.

Augustine's wish for the zealous servants, ``may it only find you wheat!'',
leaves even them unsure how the harvest will find them. In the
Postcommunion the Mass asks \latin{ut illíus salutáris capiámus efféctum},
``that we may receive the effect of that salvation,'' and since the parable's
end for the wheat is the barn, that is a prayer to be found wheat at the
harvest.

\subsection{What the cockle is, and where the field lies}
\label{sec:patience-limits}

The parable's reason applies to each of the Fathers' names for the cockle,
heretics, evil Christians, the wicked generally; but what Chrysostom forbids
is the killing of heretics, and Augustine's charge, ``let the bad change
themselves,'' is addressed to sinners among the faithful.

The servants are named differently too. Jerome says whom they stand for:
\latin{Servos patrisfamilias, ne alios accipias quam angelos}, ``by the
servants of the householder understand none other than angels'' (PL~26
col.~93). Augustine likens himself and his hearers to them. Aquinas, noting
that the reapers are angels and not servants, calls the servants
\latin{homines boni}, good men (\work{Super Mt.}\ c.~13). Honorius of Autun,
in a sermon on the parable, calls them \latin{praelati Ecclesiae qui cupiunt
malos a bonis segregare}, prelates who wish to separate the bad from the good
(\work{Speculum Ecclesiae}, the sermon headed \latin{Dominica tercia},
PL~172). Jerome's own warning to the hearer comes from \latin{ne forte}, not
from the servants.

The Lord's own gloss says ``the field is the world'' (Mt 13:38), and the
Fathers place the cockle among the faithful. Augustine holds both together in
his letter to Vincentius: the Church sown by the Son of man grows among the
tares, ``for the field is the world, and the harvest is the end of time''
(\work{Ep.}~93.31). Honorius, in the same sermon, says the mixture is allowed
\latin{ut boni per malos exerceantur et pro patientia habundantius remunerari
mereantur}, ``that the good may be exercised by the bad and may deserve to be
rewarded more abundantly for their patience.''

\begin{fourSenses}
\item[Literal.] A householder forbids his servants to weed the darnel from his
wheat before harvest, lest they uproot the wheat with it, and leaves the
separation to the reapers at the harvest; Paul tells the Colossians to bear
with and forgive one another as the Lord forgave them, and above all to have
charity.
\item[Allegorical.] Christ sows the good seed and will judge through his
angels, ``who cannot make mistakes'' (Augustine); the Church on pilgrimage
lives in ``a certain distinction without separation'' (Augustine), with wheat
and cockle in her high seats as among her laity (Augustine), until the end of
the world (Mt 13:39).
\item[Moral.] Judge slowly where the case is doubtful, because we may be
deceived (Jerome, Augustine); seek the conversion of the erring rather than
their destruction, since today's cockle may be tomorrow's wheat (Augustine,
Chrysostom); bear injuries without revenge and keep patience to the end of
life (Rabanus); bear with, forgive and love the offender after Christ's
pattern (Chrysostom and Aquinas on Col 3); pray for enemies (Jerome on Jer
29:7); and forgive in order to be heard (Bede).
\item[Anagogical.] At the harvest the angels, who cannot err, make the
separation (Augustine, \work{Sermo}~73.4); ``In the field they may be with
you, but they will not be so in the barn'' (73.1); the householder who bore
long with the cockle repays each by justice (Rabanus); and the wheat is
gathered into the barn, where, in the Lord's explanation, ``the just shall
shine as the sun, in the kingdom of their Father'' (Mt 13:43).
\end{fourSenses}
